% Глава: 13-c-compilation
% Статус: наполнено (первый прогон).
\chapter{Язык C и компиляция}
\label{ch:13-c-compilation}

\section{Язык C: зачем и что важно}

C --- низкий, близкий к машинному коду язык: предсказуемые
свойства, отсутствие скрытых механизмов, прямой доступ к
памяти. Именно поэтому ядро Linux и большинство системных
инструментов написаны на C. Для курса важно понимать:
как работает компилятор, как линковщик собирает программу,
как устроена память.

Компилируем на сервере: \cmd{gcc} запускается на вашей ВМ
\fpath{pglab.nsu.ru}, и все артефакты сборки тоже остаются на ней.

\section{Компилятор и интерпретатор}

\textbf{Компилятор} (C) переводит весь исходный код в
машинный код \emph{до} запуска; результат --- исполняемый
файл. \textbf{Интерпретатор} (Python, Ruby, большинство
JS) выполняет исходный код «на лету». Компиляция даёт
быстроту выполнения, интерпретация --- простоту запуска.

\section{Этапы компиляции C}

Командой \cmd{gcc} вы обычно делаете всё сразу, но внутри
несколько этапов:

\begin{figure}[htb]
  \centering
  \begin{tikzpicture}[node distance=1.3cm]
    \node[data] (src) {source.c};
    \node[process, right=of src] (pp) {preprocess\\[0.2em] (gcc -E)};
    \node[process, right=of pp] (asm) {assemble\\[0.2em] to .o};
    \node[process, right=of asm] (lnk) {link\\[0.2em] (ld)};
    \node[data, right=of lnk] (bin) {binary};
    \draw[flow] (src) -- (pp);
    \draw[flow] (pp) -- (asm);
    \draw[flow] (asm) -- (lnk);
    \draw[flow] (lnk) -- (bin);
  \end{tikzpicture}
  \caption{Этапы: препроцессор, компиляция в объект, линковка.}
  \label{fig:c-pipeline}
\end{figure}

\begin{itemize}
  \item \textbf{Preprocessing} (\cmd{gcc -E}) --- раскрытие
    \cmd{\#include}, \cmd{\#define}, макросов;
  \item \textbf{Compilation} --- перевод в assembly;
  \item \textbf{Assembly} --- перевод assembly в машинный
    код объектного файла \fpath{.o};
  \item \textbf{Linking} --- сборка объектных файлов и
    библиотек в один исполняемый файл.
\end{itemize}

\section{Типовой сценарий}

\begin{minted}[bgcolor=gray!6]{c}
// hello.c
#include <stdio.h>

int main(void) {
    int n = 42;
    printf("answer: %d\n", n);
    return 0;
}
\end{minted}

\begin{minted}[bgcolor=gray!6]{bash}
gcc -Wall -Wextra -o hello hello.c   # за один раз: исходник -> бинарник
./hello                              # запустить
gcc -c hello.c                       # только компиляция: hello.c -> hello.o
ld hello.o -o hello -lc              # собрать вручную, нужна библиотека
\end{minted}

\section{Флаги \cmd{gcc}}

\begin{center}
\begin{tabular}{lp{0.65\linewidth}}
\toprule
\textbf{Флаг} & \textbf{Назначение} \\
\midrule
\cmd{-c}        & только компиляция в объектный файл, без линковки \\
\cmd{-o name}   & имя выходного файла \\
\cmd{-I dir}    & каталог для поиска заголовков (\cmd{\#include}) \\
\cmd{-L dir}    & каталог для поиска библиотек на этапе линковки \\
\cmd{-l lib}    & подключить библиотеку (например \cmd{-lm} --- math) \\
\cmd{-Wall} \cmd{-Wextra} & включение большинства предупреждений \\
\cmd{-g}        & отладочная информация (для gdb, valgrind) \\
\cmd{-O0 -O1 -O2 -O3} & уровни оптимизации \\
\cmd{-std=c11}  & выбрать стандарт языка \\
\bottomrule
\end{tabular}
\end{center}

\begin{tipbox}
Всегда компилируйте с \cmd{-Wall -Wextra}: предупреждения
компилятора ловят типовые ошибки (неиспользуемые переменные,
потерянные значения) раньше, чем рантайм.
\end{tipbox}

\section{Заголовочные файлы и include-пути}

C делит код на \textbf{модули} (\fpath{.c} файлы) с
публичным интерфейсом в \fpath{.h} заголовках. \cmd{\#include
"file.h"} ищет в текущем каталоге, \cmd{\#include <file.h>} ---
в системных каталогах. Если заголовки лежат не в
системных каталогах, укажите их через \cmd{-I}.

Типовой проект:
\begin{minted}[bgcolor=gray!6]{bash}
project/
  src/main.c
  src/util.c
  src/util.h
  Makefile
\end{minted}

\begin{minted}[bgcolor=gray!6]{bash}
gcc -Wall -Wextra -Isrc -o app src/main.c src/util.c
\end{minted}

\section{Библиотеки}

\textbf{Статическая} (\fpath{libfoo.a}) --- код библиотеки
встраивается в ваш бинарник; бинарник самодостаточен.
\textbf{Динамическая} (\fpath{libfoo.so}) --- общий код
на диске; бинарник ссылается на него во время выполнения.
Динамические меньше по размеру и позволяют обновлять
библиотеку без пересборки.

\begin{minted}[bgcolor=gray!6]{bash}
gcc -o prog main.c -L./lib -lfoo      # линковка с libfoo
ldd prog                              # показать используемые общие библиотеки
\end{minted}

\section{Типичные ошибки}
\begin{itemize}
  \item \emph{undefined reference to ...} --- линковщик
    не нашёл функцию; проверьте, все ли \fpath{.c}
    переданы и не забыта ли библиотека.
  \item \emph{implicit declaration of function} --- забыли
    \cmd{\#include} или опечатались в имени; при
    \cmd{-Wall} это предупреждение, но лучше править.
  \item Выходной файл не запускается --- не тот процессор
    или отсутствует библиотека; \cmd{ldd} покажет, что
    не хватает.
  \item Код «работает, но странно» --- чаще всего это
    обращение к неинициализированной памяти; см.
    \nameref{ch:17-debugging}.
\end{itemize}

\section{Связь с остальным курсом}

Сборка проекта через \cmd{make} --- \nameref{ch:14-make};
тестирование и оптимизация --- \nameref{ch:15-testing-optimization};
память и как её распределяет ОС ---
\nameref{ch:16-virtual-memory}; отладка и поиск утечек ---
\nameref{ch:17-debugging}. Docker-образы с C-проектами
часто включают \cmd{gcc} и \cmd{make}
(\nameref{ch:20-dockerfile}).

\section{Задачи для самопроверки}
\begin{enumerate}
  \item Объясните четыре этапа компиляции C. Что
    генерирует каждый из них?
  \item Чем \cmd{gcc -c} отличается от \cmd{gcc}?
  \item Как скомпилировать проект из нескольких
    \fpath{.c} с заголовками в подкаталоге?
  \item В чём разница между статической и динамической
    библиотекой?
  \item Какие флаги вы бы использовали для отладочной
    и для оптимизированной сборки?
\end{enumerate}
